% ---------------------------------------------------------------
%  PHAN III -- KET LUAN, KIEN NGHI + TAI LIEU + PHU LUC
% ---------------------------------------------------------------
\section{KẾT LUẬN VÀ KIẾN NGHỊ}

\subsection{Kết luận}

Sáng kiến đã xây dựng và đưa vào sử dụng thực tế một hệ thống ngân hàng đề môn
Toán trên nền web dùng chung cho cả giáo viên và học sinh, với ba kết quả
chính:

\begin{enumerate}[label=(\arabic*), leftmargin=1.2cm, itemsep=3pt]
  \item \textbf{Về mặt giải pháp:} đã chứng minh được một hướng đi khác với hai
        cách làm phổ biến hiện nay. Không dùng ngân hàng câu hỏi tĩnh --- vì
        tĩnh thì dùng một lần là hết; cũng không giao cho trí tuệ nhân tạo sinh
        đề --- vì không kiểm chứng được. Thay vào đó, câu hỏi do chính giáo
        viên viết ra dưới dạng hàm sinh bằng Python, nên vừa tự sinh được vô
        hạn biến thể, vừa bảo đảm tuyệt đối không sai đáp án và không vượt
        khung chương trình.
  \item \textbf{Về mặt sản phẩm:} hệ thống đang chạy thật tại địa chỉ
        \url{https://nganhangdechv.tech}, thực hiện trọn vòng từ một câu yêu
        cầu tiếng Việt đến đề PDF, lời giải, bài làm trực tuyến, điểm số, sổ
        điểm Google Classroom và bảng thống kê của giáo viên.
  \item \textbf{Về mặt chuyên môn:} các quy ước sư phạm của người viết --- mã
        định danh gắn với yêu cầu cần đạt, một câu Đúng/Sai gồm bốn ý ở bốn mức
        độ, chia câu theo tỉ lệ số tiết, mỗi bài nhiều nhất một câu vận dụng
        cao --- đã được mã hoá thành thuật toán, nghĩa là chúng được thực hiện
        chính xác và nhất quán ở mọi đề, điều mà làm thủ công rất khó bảo đảm.
\end{enumerate}

Bài học lớn nhất tôi rút ra không thuộc về kĩ thuật: \textbf{công cụ chỉ mạnh
khi người giáo viên có quy ước chuyên môn rõ ràng trước khi ngồi vào máy}. Mọi
lần hệ thống cho kết quả sai trong quá trình xây dựng đều quy về một trong hai
nguyên nhân --- hoặc tôi chưa nói rõ quy ước, hoặc tôi đã giao cho máy phần
việc lẽ ra phải do người quyết định.

\subsection{Kiến nghị}

\textbf{Với tổ chuyên môn.} Thống nhất một bộ ma trận chung cho từng loại bài
kiểm tra và phân công mỗi giáo viên viết hàm sinh câu hỏi cho một vài bài mình
dạy kĩ nhất. Nếu làm được, ngân hàng chung của tổ sẽ lớn rất nhanh và mỗi người
đều được hưởng thành quả chung.

\textbf{Với nhà trường.} Tạo điều kiện về phòng máy hoặc cho phép học sinh dùng
điện thoại trong tiết luyện tập có kiểm soát, để việc làm bài trực tuyến diễn
ra ngay tại lớp; đồng thời hỗ trợ tổ chức một buổi tập huấn nội bộ để đồng
nghiệp làm quen với hệ thống.

\textbf{Với Sở Giáo dục và Đào tạo.} Tôi kiến nghị có cơ chế chia sẻ ngân hàng
câu hỏi dạng hàm sinh giữa các trường trong tỉnh, kèm một chuẩn mã định danh
thống nhất. Khi ngân hàng đủ lớn, mọi giáo viên trong tỉnh đều có thể lấy được
đề đúng ma trận cho lớp mình --- đó là phần giá trị vượt xa phạm vi một trường
mà tôi mong sáng kiến này góp được.

\bigskip
\noindent Tôi xin cam đoan sáng kiến này là của bản thân tôi thực hiện, không
sao chép của người khác. Trong quá trình thực hiện tôi có sử dụng trợ lí trí
tuệ nhân tạo để hỗ trợ viết và gỡ lỗi mã nguồn, điều này đã được nêu rõ tại Mục
II.3; toàn bộ quy ước chuyên môn, nội dung toán học và các quyết định sư phạm
là do tôi đưa ra và chịu trách nhiệm.

\vspace{1cm}
\begin{flushright}
\begin{minipage}{6cm}
\centering
\textit{\DIADIEM, ngày \dots\ tháng \dots\ năm 2026}\\[2pt]
\textbf{Người viết sáng kiến}\\
\textit{(Ký và ghi rõ họ tên)}\\[2.5cm]
\TACGIA
\end{minipage}
\end{flushright}

\clearpage

% ---------------------------------------------------------------
\section*{TÀI LIỆU THAM KHẢO}
\addcontentsline{toc}{section}{TÀI LIỆU THAM KHẢO}
\setcounter{secnumdepth}{0}

\begin{enumerate}[leftmargin=1.2cm, itemsep=3pt]
  \item Bộ Giáo dục và Đào tạo (2018), \textit{Chương trình Giáo dục phổ thông
        --- Chương trình môn Toán}, ban hành kèm theo Thông tư số
        32/2018/TT-BGDĐT ngày 26/12/2018.
  \item Bộ Giáo dục và Đào tạo (2021), \textit{Thông tư số 22/2021/TT-BGDĐT
        ngày 20/7/2021 quy định về đánh giá học sinh trung học cơ sở và học
        sinh trung học phổ thông}.
  \item Bộ Giáo dục và Đào tạo (2024), \textit{Quyết định số 764/QĐ-BGDĐT ngày
        08/3/2024 về việc phê duyệt cấu trúc định dạng đề thi Kỳ thi tốt nghiệp
        trung học phổ thông từ năm 2025}.
  \item Thủ tướng Chính phủ (2022), \textit{Quyết định số 131/QĐ-TTg ngày
        25/01/2022 phê duyệt Đề án ``Tăng cường ứng dụng công nghệ thông tin và
        chuyển đổi số trong giáo dục và đào tạo giai đoạn 2022--2025, định
        hướng đến năm 2030''}.
  \item Bộ Chính trị (2024), \textit{Nghị quyết số 57-NQ/TW ngày 22/12/2024 về
        đột phá phát triển khoa học, công nghệ, đổi mới sáng tạo và chuyển đổi
        số quốc gia}.
  \item Chính phủ (2012), \textit{Nghị định số 13/2012/NĐ-CP ngày 02/3/2012 ban
        hành Điều lệ Sáng kiến}.
  \item Nhóm tác giả (2022), \textit{Toán 10 --- Bộ sách Kết nối tri thức với
        cuộc sống}, Nhà xuất bản Giáo dục Việt Nam.
  \item Tài liệu kĩ thuật của hệ thống Ngân hàng đề AI (sổ tay kĩ thuật nội bộ,
        20 tài liệu, cập nhật đến tháng 9/2026).
\end{enumerate}

\ghichu{Bổ sung thêm sách, bài báo về đo lường và đánh giá mà chị thực sự có
đọc. Không nên liệt kê tài liệu chưa đọc --- hội đồng có thể hỏi.}

\clearpage

% ---------------------------------------------------------------
\section*{PHỤ LỤC}
\addcontentsline{toc}{section}{PHỤ LỤC}

\textbf{Phụ lục 1.} Phiếu khảo sát giáo viên và phiếu khảo sát học sinh.

\textbf{Phụ lục 2.} Một đề kiểm tra định kì hoàn chỉnh do hệ thống sinh ra, kèm
bản lời giải.

\textbf{Phụ lục 3.} Hai mã đề cùng cấu trúc, khác số liệu (minh chứng cho Giải
pháp 5).

\textbf{Phụ lục 4.} Bảng ma trận chi tiết của đề ở Phụ lục 2.

\textbf{Phụ lục 5.} Tài khoản dùng thử dành cho hội đồng chấm:
\url{https://nganhangdechv.tech} --- tên đăng nhập \dotfill, mật khẩu \dotfill.

\ghichu{Phụ lục 5 rất nên có: tạo sẵn một tài khoản học sinh dùng thử và một
tài khoản giáo viên để hội đồng tự vào bấm. Đây là điểm mạnh mà các sáng kiến
khác không có --- sản phẩm chạy được thật, kiểm chứng được ngay.}

\ghichu{In kèm Phụ lục 2 và 3 ở dạng bản in thật (không phải ảnh chụp màn
hình): hai mã đề đặt cạnh nhau để thấy rõ cùng dạng toán, khác số liệu.}

